\section{A Story of Love and Death}

\paragraph{In Memoriam to my mother and father}


\textbf{Update}: Three years ago today, my father entertained family guests and had a stroke an hour later. Two months later, both my mother and my father had died, one from disease, the other from loneliness. This is the story of their last days together.

My father had a severe stroke in a year ago today; he had been taking care of my mother with Alzheimer's. Really, neither of us (i.e., my sister and I) knew the extent of that care. He seldom complained, but we came to see how much care she needed: from showering, to getting dressed, to eating. She couldn't have been good company for him because she couldn't even follow a TV show. Every time I visited, I would go over family photographs with her. She knew she had a family, but she couldn't place the names and faces. It was extremely frustrating to her.

Yet, they didn't want to move to assisted living or even to a spare room at my sister's house --- they valued their independence right up until the end. My father had retired when he was 50 and he and my mother have been inseparable for the past 37 years. Just before his stroke, he had been entertaining family that afternoon just as he loved to do. Friends, family, food, drink ... that was how he preferred to live his life.

In her state of mind, my mother never did adjust after my father's stroke, even though my sister took her in and gave her good care. Several weeks later, she had a mild problem (intestinal blockage) --- in retrospect, the doctor may have overreacted and sent her to the hospital, where they took their time tending to her. One night, she tried to get out of bed, fell down, and broke her hip. That's how she ended up in the same rehab center where my father was staying. My mother's Alzheimer was getting progressively worse there.

Yet, she recognized my father. She would wheel her chair right up to him, caress his face, hold hands, proclaim her love, and tell him to get better. It was so touching, the center staff would gather around just to witness it.

My father passed away peacefully around 7:35 on a Monday night in June --- I was the only family member there.

Right afterwards my mother (who is at the same rehab center in a room directly across the hall), woke up from a deep sleep and started crying out ``Chuck Chuck''. Any other time you could have asked her for her husband's name, but the queen would have guessed Rumpelstiltskin's name before my mother would have guessed ``Chuck''. But this night, the name was spoken with confidence and authority.

I went to her to give her a hug and she embraced and kissed me, over and over, repeatedly calling ``Chuck''.

Her touches and kisses were full of passion ... and sadness. She refused to lighten her grip on me. ``Don't leave me'' and ``I love you, baby'', she said, over and over. The intensity was both unexpected and overpowering as it gave me a glimpse into their life together, a part that they kept hidden from the children. It had to have come from a deep and primordial part of her mind that, once freed of its inhibitions, displayed itself in all its raw emotion.

I'm sure there was an intense spiritual bond between the two of them at that moment.

At some deep level she became aware that my father --- her husband --- had died, without anyone having told her. She was never the same and was constantly crying. The last conversation I had with her was a week later. We were sitting outside at the rehab center when she suddenly got one of her lucid moments. She looked at me and asked ``Papa died, didn't he?'' I had to answer truthfully. With that she started crying. I was never able to hold another conversation with her, and within 10 days, as her ultimate act of devotion she passed away herself, dying from heartache and grief.

The final end of the story only came few months later. It was common for me to dream, but this one was particularly vivid.

\begin{quotationx}
I was with Nana and Papa at a large arena for a sporting event ... we were waiting in the lobby holding tickets for the lower level at cheap seats.

We decided to exchange them for \$50 seats on the upper level. I went with Nana upstairs. There I ran into my two sons who had already made the exchange. They had seats numbered 4, 5 and 6. Nana stood behind them, silent and ethereal, with her arms outstretched.

I went downstairs to fetch Papa. Though he looked pale and shuffled slowly as we walked, I thought to myself: ``He doesn't seem so bad considering he just had a stroke.''

Then I realized that there were only three seats ... for me and the two boys and my mother was already there. There was no place there for my father. I slowly led him up ... 
\end{quotationx}



\flrightit{Posted on 2008-04-21 by Cologero}

\begin{center}* * *\end{center}

\begin{footnotesize}
\begin{sffamily}
\texttt{Anna Cue on 2008-04-26 at 03:22 said:}

I enjoyed reading it.

\hfill

\texttt{Katerina on 2010-04-18 at 17:52 said:}

hauntingly beautiful... 

\hfill

\texttt{cologero on 2010-04-18 at 19:22 said:}

Thanks to \url{http://www.freecrosses.com/celtic\_cross.htm} for celtic cross image.

\hfill

\texttt{Barb on 2010-04-21 at 09:19 said:}

I once had a relationship as deep as the one you speak of... ..death and departure hurt. Fortunately, we're left with the memories -- hopefully holding onto the good ones. I'm sure your parents went through trials... and as I knew them... .they truly loved each other and help together. Having known your parents I used to admire their love of each other... .it's a rare thing ``these days.''

I too have had dream in which your parents have had many conversations with me... .especially your father. I was thrilled to know them... they were both wonderful... .and your sister, true to form was very giving.

You are in my prayers.

Barb

\hfill

\texttt{Cologero on 2011-10-30 at 07:30 said:}

From a recent dream ... 

As I was driving, I got a call from my father. He invited me to dinner, promising to make a special dish. Believing I found a shortcut, I took a right turn on an unknown road, assuming it would take me straightaway to his house. Instead, the road led to a cul de sac. I got out of the car to look. There was a steep drop off, but I was struck by the scene. I recall thinking that it was the most beautiful view I had ever seen in my life. The colors were vivid, there were flowers and a small compound of elegant cottages in the distance. I clambered down for a closer look. However, I was unable to reach what I was seeing.

Returning involved climbing back up the cliff, and I thought that I might not be able to get back. As I approached the small ledge at the top, there were two older women. I asked about the place. They told me it was a very exclusive neighborhood with strict zoning regulations, and we discussed how beautiful it was. At the same time, I noticed a younger woman also climbing up onto the ledge. She had been visiting the scene, also, and told me she had little success there ... I don't recall what she was doing, perhaps selling something.

\hfill

\texttt{Charlotte Cowell on 2011-10-30 at 11:50 said:}

You're definitely tuned in today C! I was just about to think on this when I turned on the computer today, having promised to more than once :-)

The first part seems to me an uncomplicated foresight into the heavenly abode you'll share with your parents, I think this is probably obvious? Also, perhaps, a reminder, to always look beyond, to leave no stone unturned in the quest for beauty, love and enlightenment.

Part two is unusual in the sense that you found yourself at the bottom of the valley without having knowingly climbed down, unless you just didn't' remember that? (in lucid states like this time and space are irrelevant, so if your will directs you to a certain place you will simply arrive there in no time at all, in a spontaneous and instantaneous arrival at where you are `supposed' to be').

I can tell you definitely that climbing up in dreams is good, even if it's a struggle, and climbs up hills and cliffs generally are in dreams as well as real life. They are straightforward representations of the tasks and struggles most pertinent to our life mission and quest, the fact you went upwards was a sign you were making progress.

Not sure who any of the women, older or younger are, the may represent people you already know, but it's telling that you all agreed on how lovely it was -- a positive place, then, do you remember at all how you felt towards both the place and them? (by the way the mention of strict zoning and exclusivity may have been a test or something to do with ego, I'm not sure).

The younger woman may also be someone you know, or (as might the others) an elemental being of the kind that frequent dreams and offer signs, tests, pointers, etc. Was she telling you that the place as a red herring of sorts (too exclusive, therefore not universal?), or was it simply the case that this was a place for you and your parents and `not for' anyone else, but your private refuge with them?

It may have been a faerie place, which I suggest because of the emphasis on natural beauty.

The impression of very vivid colours is a spiritual sign in dreams and a measure of the lucidity -- I was taught once that there is a `plane of colour' and it's always worth noting any colours that stand out in the dream state, and also any sounds as so much that occurs there is in silence.

But you may of course disagree with all of this and already know exactly what it means... Cx

\hfill

\texttt{Cologero on 2011-10-30 at 12:57 said:}

I took the right turn because my father was not at the place where I would have expected him; compare this to the dream about the sporting event when he was not at the upper level seats. His time of purgation seems to have been completed.

The older women represent Moira, the Fates. The third fate is replaced by the young woman, the \textit{Divine Sophia}\footnote{\url{http://www.gornahoor.net/?p=387}} who has appeared to me in many dreams in many ways.

\hfill

\texttt{Charlotte Cowell on 2011-10-30 at 13:48 said:}

nice (the earlier Divine Sophia dream I mean, that sounds amazing)! Anyways, I'm glad your father is doing well and has passed over in a good way.

\hfill

\texttt{Charlotte Cowell on 2011-11-01 at 20:44 said:}

I just heard this... .

\url{http://www.youtube.com/watch?v=65Lxd2kc8Rs&feature=related}


Cx

\end{sffamily}
\end{footnotesize}

